\appendix
\section{Selection Laws And Bounded Theorems}

\subsection{The \texttt{P2R.B*} Selection Law}

Let $T = \texttt{P2R.B*}$. Define $F_T(s)=1$ when direct local source/IR
lowering emits $T$, and define $G_T(b)=1$ when cubin-side substitution
materializes runnable $T$.

\begin{theorem}
On the current SM89 stack, the observed \texttt{P2R.B*} gap is a form-selection
gap rather than an opcode-existence gap.
\end{theorem}

\begin{proof}
The repository evidence establishes three facts: (i) direct source/IR search
does not emit \texttt{P2R.B1}, \texttt{P2R.B2}, or \texttt{P2R.B3}; (ii) plain
\texttt{P2R} neighborhoods are reproduced locally; and (iii) cubin-side
substitution materializes runnable local \texttt{P2R.B1/B2/B3}. Therefore the
missing condition lies in opcode-form selection under the tested lowering path,
not in local opcode invalidity.
\end{proof}

\subsection{The \texttt{UPLOP3} Structural Law}

\begin{theorem}
For the tested local cubin contexts, \texttt{ULOP3 -> UPLOP3} is structurally
invalid while \texttt{PLOP3 -> UPLOP3} is structurally valid.
\end{theorem}

\begin{proof}
The invalid substitution class decodes as illegal under local disassembly,
whereas the predicate-logic substitution class decodes as real
\texttt{UPLOP3.LUT} and launches successfully in multiple local contexts.
\end{proof}

\subsection{The \texttt{UPLOP3} Runtime-Class Law}

\begin{theorem}
The tested local \texttt{UPLOP3} substitutions partition into inert and
stable-but-different runtime classes.
\end{theorem}

\begin{proof}
Differential fuzzing separates contexts whose outputs match baseline across all
tested patterns from contexts whose outputs diverge on at least one tested
pattern, while compute-sanitizer remains clean. This yields the inert versus
stable-but-different partition used throughout the current Ada-only frontier.
\end{proof}
